Systems of twinned systems (SoTS) represent the convergence of system of systems (SoS) and digital twins (DTs). While SoS research emphasizes the orchestration of independent, heterogeneous systems to achieve collective goals \cite{maier1998architecting}, DTs focus on establishing a synchronized digital representation of physical assets. This enables real-time monitoring, analysis, and control \cite{kritzinger2018digital, tao2018digital}. Integrating these ideas creates opportunities to extend both. DTs gain scalability and coordination across multiple interacting systems, while SoS benefit from intelligent, data-driven representations of their constituents. To explore this convergence, a systematic literature review (SLR) was conducted to identify, classify, and analyze combining DT and SoS. Our goal was to understand the characteristics of SoTS, their components, and constituent systems, as well as to identify the key limitations, challenges, and research opportunities in the field.


\section{Research Questions}\label{sec:literature-review-rqs}
This review aimed to achieve answer the following objectives:

\vspace{0.5em}
\noindent \textbf{\large Conceptual Motivation and Integration}
\vspace{-1em}

\begin{itemize}
    \item \rquestion{RQ1. Why integrate Digital Twins into System of Systems?} This question seeks to uncover the motivations and intents behind forming SoTS, the domains where they are applied, and the key challenges that arise when multiple DTs are organized as a coordinated system.
    \item \rquestion{RQ2. How are Digital Twins integrated into System of Systems?} This question examines the architectural principles used to organize DTs within an SoS. We analyze the nature of constituent units (physical, digital, or hybrid) and classify the resulting SoS structures.
\end{itemize}
    Together, these questions reveal not only why SoTS are pursued, but also how they are practically realized in current research.

\vspace{0.5em}
\noindent \textbf{\large Technical Characterization}
\vspace{-1em}
\begin{itemize}
    \item \rquestion{RQ3. What are the technical characteristics of Digital Twins in SoTS?} This question analyzes how DTs function as constituents within SoTS — including their autonomy, supported services, modelling formalisms, and interaction capabilities.
    \item \rquestion{RQ4. What are the technical characteristics of System of Systems in SoTS?} This question examines SoS-level properties in SoTS, such as SoS classifications, coordination mechanisms, and the types of emergent behaviour addressed.
\end{itemize}
Together, these questions provide a detailed view of how both DTs and SoS contribute technically to the structure and operation of SoTS.

\vspace{0.5em}
\noindent \textbf{\large Maturity and Realization}
\vspace{-1em}
\begin{itemize}
    \item \rquestion{RQ5. How are non-functional properties addressed in SoTS?}
    This question investigates how reliability and security are treated. Whether they are simply acknowledged, incorporated architecturally, explicitly modelled, or empirically evaluated.
    \item \rquestion{RQ6. What is the level of maturity in SoTS research and technology?} This question assesses both technical maturity and research maturity, while also examining whether standards are adopted.
    \item \rquestion{RQ7. What technological choices are used to implement SoTS?}  
    This question surveys the platforms, languages, and frameworks that enable SoTS deployment in practice.
\end{itemize}
Altogether, these questions evaluate how well current SoTS research supports dependable operation, readiness for real-world deployment, and the practical implementation paths available today.
